\chapter{Точная LCF-проверка минимальной Стайнспринг-дилатации}
\label{ch:v8-minimal-covariant-stinespring-lcf-migration-gate}

\section{Одношаговый канал}

Двенадцать самосопряжённых межсекторных стрелок операторов $D_a$ проверены как
точная ортогональная рамка,
\begin{equation}
 \Tr(D_aD_b)=2\delta_{ab}.
\end{equation}
Для $G=\sum_aD_a^2$ точный спектр равен
\begin{equation}
 \Spec G=\{0^9,1^6,2^3,3^2,6\}.
 \label{eq:v8-lcf-stinespring-gram-spectrum}
\end{equation}
Поэтому выражения
\begin{equation}
 K_0(p)=\sqrt{I-pG},\qquad K_a(p)=\sqrt p\,D_a
 \label{eq:v8-lcf-stinespring-kraus-family}
\end{equation}
задают канал точно при $0\le p\le1/6$.

На рациональной контрольной точке $p=1/12$ обе полноты Краус-семейства
проверены символически. Тем самым карта полностью положительна, унитальна и
сохраняет след. Инвариантность концевой алгебры проверена на всех $221$
её матричных единицах.

\section{Минимальность и ковариантность}

Для каждого внутреннего $p>0$ векторизации $K_0,K_1,\ldots,K_{12}$ линейно
независимы. Следовательно, Краус-ранг, ранг Чой и минимальная комплексная
размерность среды совпадают и равны $13$. Среда имеет канонический вид
\begin{equation}
 \mathbb C\lvert0\rangle\oplus\mathcal E_{\rm cross}^{\mathbb C},
 \qquad \dim_{\mathbb R}\mathcal E_{\rm cross}=12.
\end{equation}
Её скачковая часть является комплексификацией уже имеющейся межсекторных стрелок рамки,
а не новым физическим мультиплетом. Калибровочная ковариантность следует из
точного ортогонального действия на этой рамке и независимости Краус-суммы
от выбора ортогонального базиса.

\section{Граница непрерывного времени}

В точке $p=0$ выполнено
\begin{equation}
 K_0(0)=I,\qquad K_0'(0)=-\frac12G,
\end{equation}
поэтому производная канала совпадает с межсекторных стрелок GKSL-генератором.
Однако конечное семейство не является полугруппой: точная матричная единица
даёт
\begin{equation}
 \Phi_{1/50}\!\circ\Phi_{3/100}\ne\Phi_{1/20}.
\end{equation}

\begin{theorem}[Минимальная одношаговая дилатация]
Репер межсекторных стрелок канонически задаёт калибровочно-ковариантный одношаговый
канал с минимальной средой размерности $13$ для $0<p<1/6$; новый физический
носитель для этого не требуется.
\end{theorem}

\begin{nogocriterion}[Одношаговый параметр не является временем]
Нельзя отождествлять $p$ с физическим временем или выведенной скоростью.
Тангенс в нуле является GKSL-генератором, но конечные карты не удовлетворяют
закону полугруппы. Непрерывный шум требует отдельного предела слабых столкновений либо
внутреннего часового механизма.
\end{nogocriterion}

\begin{protocol}
\begin{enumerate}
 \item Гильберт--Шмидт Грам рамки: \textbf{$2I_{12}$};
 \item спектр $G$: \textbf{$\{0^9,1^6,2^3,3^2,6\}$};
 \item допустимое окно: \textbf{$0\le p\le1/6$};
 \item полнота Краус-семейства: \textbf{точно};
 \item сохранение следа: \textbf{точно};
 \item единиц концевых секторов проверено: \textbf{$221$};
 \item внутренний Краус/Чой-ранг: \textbf{$13$};
 \item минимальная среда: \textbf{$13$};
 \item калибровочная ковариантность: \textbf{точно};
 \item GKSL-тангенс: \textbf{точно};
 \item конечная полугруппа: \textbf{нет}.
\end{enumerate}
\end{protocol}

Машинный сертификат сохранён в
\path{s2t_v8_minimal_covariant_stinespring_lcf_migration_gate_results.json}.